\documentclass[aspectratio=169, 11pt, xcolor={table}]{beamer}
% Shared style for the course slides, based on presentation_text/main.tex.
% Dark blue removed: every slide (cover included) uses the light background
% of the template's content slides, with teal + orange as accents.

\usepackage{amsmath}
\usepackage{amssymb}
\usepackage{booktabs}
\usepackage{colortbl}
\usepackage{tabularx}
\usepackage{tikz}
\usepackage[skins]{tcolorbox}
\usepackage{fontawesome5}
\usepackage{polyglossia}

\usetikzlibrary{positioning, calc, arrows.meta, shapes.geometric}

%------------------------------------------------
% Languages and fonts (compile from slides/ with XeLaTeX or tectonic)
%------------------------------------------------
\setdefaultlanguage[numerals=maghrib]{arabic}
\setotherlanguage{english}
\usefonttheme{serif}

\newfontfamily\arabicfont[
    Script=Arabic, Scale=0.98, Path=../fonts/, Extension=.ttf,
    UprightFont=*-Regular, BoldFont=*-Bold,
    ItalicFont=*-Italic, SlantedFont=*-Italic,
    BoldItalicFont=*-BoldItalic, BoldSlantedFont=*-BoldItalic]{Amiri}
\newfontfamily\arabicfontsf[
    Script=Arabic, Scale=0.98, Path=../fonts/, Extension=.ttf,
    UprightFont=*-Regular, BoldFont=*-Bold,
    ItalicFont=*-Italic, SlantedFont=*-Italic,
    BoldItalicFont=*-BoldItalic, BoldSlantedFont=*-BoldItalic]{Amiri}
% Heading font
\newfontfamily\khatt[Script=Arabic, Scale=1.05, Path=../fonts/,
    BoldFont=ae_AlBattar.ttf]{ae_AlBattar.ttf}

%------------------------------------------------
% Colors (template palette without the navy tones)
%------------------------------------------------
\definecolor{brandTeal}{HTML}{00B4D8}
\definecolor{brandOrange}{HTML}{FF7043}
\definecolor{brandLight}{HTML}{F8F9FA}
\definecolor{brandDark}{HTML}{263238}
\definecolor{brandMuted}{HTML}{6C757D}
\colorlet{tealDeep}{brandTeal!70!black}

\setbeamercolor{background canvas}{bg=brandLight}
\setbeamercolor{normal text}{fg=brandDark}
\setbeamercolor{structure}{fg=brandTeal}
\setbeamertemplate{navigation symbols}{}

% beamer.cls sets \raggedright globally; under RTL that centers short lines
% or pushes them left, so go back to justified text (flush right).
\rightskip=0pt\relax
% Lists call \raggedright too; make it a no-op for RTL slides.
\AtBeginDocument{\renewcommand{\raggedright}{\rightskip=0pt\relax}}

%------------------------------------------------
% Frame title, footer, lists
%------------------------------------------------
\setbeamertemplate{frametitle}{%
    \vspace{0.35cm}
    \noindent{\khatt\bfseries\Large\color{brandDark}\insertframetitle}
    \par
    \vskip-2pt
    \noindent{\color{brandTeal}\rule{\textwidth}{1.5pt}}
    \vspace{0.05cm}
}

% Footer: page number only
\setbeamertemplate{footline}{%
    \hbox to \paperwidth{\hfill
        {\color{brandMuted}\scriptsize\LR{\insertframenumber\,/\,\inserttotalframenumber}}%
        \hspace{0.6cm}}%
    \vspace{0.25cm}}

\setbeamertemplate{itemize item}{\color{brandTeal}\small\faChevronLeft}
\setbeamertemplate{itemize subitem}{\color{brandMuted}\small\faAngleLeft}
\setbeamertemplate{enumerate item}{\color{tealDeep}\bfseries\insertenumlabel)}

% Tables: every cell framed by a visible border
\newcolumntype{R}{>{\raggedleft\arraybackslash}X}
\newcolumntype{C}{>{\centering\arraybackslash}X}
\newcommand{\ras}[1]{\cellcolor{brandTeal!20}\textbf{#1}}
% Colored text inside a table cell (\leavevmode keeps it on the row baseline)
\newcommand{\cell}[2]{\leavevmode\textcolor{#1}{#2}}
\renewcommand{\arraystretch}{1.35}
\setlength{\arrayrulewidth}{0.6pt}
\arrayrulecolor{brandMuted!70}

%------------------------------------------------
% Boxes (the two boxes of the template's 2nd slide)
%------------------------------------------------
\tcbset{
    brandbox/.style={
        enhanced, arc=4pt, outer arc=4pt, boxrule=1.5pt,
        coltitle=white, coltext=brandDark, fonttitle=\khatt\bfseries,
        toptitle=0.8mm, bottomtitle=0.8mm,
        top=1.2mm, bottom=1.2mm, left=3mm, right=3mm,
        before skip=4pt, after skip=4pt}
}
% Teal box: definitions, objectives, summaries
\newtcolorbox{tbox}[2][]{brandbox, colback=brandTeal!10,
    colframe=brandTeal, title={#2}, #1}
% Orange box: warnings
\newtcolorbox{obox}[2][]{brandbox, colback=brandOrange!10,
    colframe=brandOrange, title={#2}, #1}
% Card (template's 3rd slide, in teal)
\newtcolorbox{card}[2][]{enhanced, colback=brandTeal!6, colframe=brandTeal,
    arc=6pt, boxrule=1.2pt, halign=center, halign title=center,
    coltitle=white, coltext=brandDark, fonttitle=\khatt\bfseries,
    top=1.5mm, bottom=1.5mm, left=2mm, right=2mm, title={#2}, #1}
% Prompt text written to the tool
\newtcolorbox{amr}[1][brandTeal]{enhanced, arc=1mm, boxrule=0pt,
    colback=white, overlay={\fill[#1] (frame.north east) rectangle ([xshift=-2.5pt]frame.south east);},
    left=6pt, right=8pt, top=2pt, bottom=2pt,
    before skip=3pt, after skip=3pt}

%------------------------------------------------
% TikZ (Arabic right-to-left inside nodes)
%------------------------------------------------
\newcommand{\nodecolor}{}
\tikzset{
    every node/.append style={
        execute at begin node={\nodecolor\everypar{\beginR}\beginR},
        execute at end node=\endR},
    lawn/.code={\def\nodecolor{\color{#1}}},
    khatwa/.style={draw=#1, fill=#1!12, very thick, rounded corners=5pt,
        minimum height=0.95cm, text width=1.9cm, align=center,
        font=\khatt\bfseries\small},
    khatwa/.default=brandTeal,
    sahm/.style={-{Stealth[length=2.5mm]}, thick, color=brandMuted},
    rang/.style={draw=#1, fill=#1!12},
    raqm/.style={circle, fill=#1, lawn=white, inner sep=1.5pt,
        font=\small\bfseries},
}

% Horizontal chain of steps, right to left:
% \flow[width]{{a}, {b\\c}, {d}}
\newcommand{\flow}[2][1.9cm]{%
    \begin{tikzpicture}[node distance=5mm]
        \foreach \t [count=\i] in {#2} {
            \ifnum\i=1
                \node[khatwa, text width=#1] (f\i) {\t};
            \else
                \pgfmathtruncatemacro\p{\i-1}
                \node[khatwa, text width=#1, left=of f\p] (f\i) {\t};
                \draw[sahm] (f\p) -- (f\i);
            \fi
        }
    \end{tikzpicture}}

%------------------------------------------------
% Cover (template's 1st slide on the light background)
%------------------------------------------------
\newcommand{\cover}[3]{%
\begin{frame}[plain]
    \begin{tikzpicture}[remember picture, overlay]
        \draw[brandTeal, line width=1.5pt, opacity=0.3]
            (current page.south east) ++(-1,1) circle (4cm);
        \draw[brandTeal, line width=0.8pt, opacity=0.25]
            (current page.north west) ++(2,2) circle (6cm);
        \draw[brandOrange, line width=0.5pt, opacity=0.25]
            (current page.north west) ++(2,2) circle (7cm);
        \node[anchor=north east] at ($(current page.north east)+(-1.0,-0.8)$) {
            \begin{tikzpicture}[scale=0.35]
                \fill[brandTeal] (0,0) -- (0.8,1.2) -- (1.6,0) -- cycle;
                \fill[brandOrange] (0.8,1.2) -- (1.6,2.4) -- (2.4,1.2) -- cycle;
            \end{tikzpicture}
        };
    \end{tikzpicture}

    \vspace{-0.05cm}
    {\hspace*{1.1cm}\khatt\bfseries\color{brandDark}
        المدرسة العليا للأساتذة ببوزريعة ــ ملحقة بومرداس\par}

    \vspace{1.2cm}
    \begin{minipage}{0.85\textwidth}
        {\color{tealDeep}\khatt\bfseries #1} \\[0.3cm]
        {\raggedleft\color{brandDark}\Huge\khatt\bfseries #2\par} \vspace{0.3cm}
        {\color{brandMuted}\large #3}
    \end{minipage}

    \vspace{1.2cm}
    \begin{minipage}{0.85\textwidth}
        \color{brandMuted}\small
        \begin{tabular}{@{}ll}
            \textbf{الأستاذ:} & د. بن عبد الرزاق يوسف \\
            \textbf{الفئة:}   & طلبة قسم اللغة العربية \\
        \end{tabular}
    \end{minipage}
\end{frame}}

\title{كيف يتعلم الذكاء الاصطناعي؟}

\begin{document}

\cover{درس الفصل الثاني}
{كيف يتعلم الذكاء الاصطناعي؟}
{البيانات ــ التدريب ــ الاختبار}

\begin{frame}{أهداف الدرس}
    \begin{tbox}{\faBullseye\quad أهداف الدرس}
        في نهاية هذا الدرس، تستطيع أن:
        \begin{itemize}
            \item تشرح مراحل تعلم الآلة: من البيانات إلى النتيجة.
            \item تفرّق بين بيانات التدريب وبيانات الاختبار.
            \item تفهم لماذا يجب أن تكون البيانات متنوعة.
            \item تتعرف على الانحياز الخوارزمي وطرق التقليل منه.
        \end{itemize}
    \end{tbox}
\end{frame}

\begin{frame}{التلميذ والآلة: تشبيه بسيط}
    يتعلم التلميذ الإعراب من أمثلة كثيرة مُعرَبة، ثم يُختبر بجمل جديدة.
    والآلة تتعلم بطريقة قريبة: تكتشف في الأمثلة \textbf{الأنماط}، أي الأشياء التي تتكرر.
    \begin{center}
    \begin{tikzpicture}[
        t/.style={thick, rounded corners=4pt,
            text width=4.2cm, align=center, minimum height=0.75cm}]
        \node[font=\khatt\bfseries, lawn=brandMuted] at (3,0.75) {التلميذ};
        \node[font=\khatt\bfseries, lawn=tealDeep] at (-3,0.75) {الآلة};
        \node[t, rang=brandMuted] (t1) at (3,0)    {يقرأ أمثلة مُعرَبة};
        \node[t, rang=brandMuted] (t2) at (3,-1.15) {يلاحظ القاعدة};
        \node[t, rang=brandMuted] (t3) at (3,-2.3) {يُعرب جملًا جديدة};
        \node[t, rang=brandTeal] (a1) at (-3,0)    {تأخذ بيانات التدريب};
        \node[t, rang=brandTeal] (a2) at (-3,-1.15) {تكتشف الأنماط};
        \node[t, rang=brandTeal] (a3) at (-3,-2.3) {تُختبر ببيانات جديدة};
        \draw[sahm] (t1)--(t2); \draw[sahm] (t2)--(t3);
        \draw[sahm] (a1)--(a2); \draw[sahm] (a2)--(a3);
        \foreach \i in {1,2,3}
            \draw[dashed, thick, brandMuted!60] (t\i) -- (a\i);
    \end{tikzpicture}
    \end{center}
\end{frame}

\begin{frame}{مراحل التعلم}
    يمر تعلم الآلة بخمس مراحل متتابعة:
    \begin{center}
        \flow[1.8cm]{{البيانات}, {التدريب}, {النموذج}, {الاختبار}, {النتيجة}}
    \end{center}
    \small
    \begin{tabularx}{\textwidth}{|>{\centering\arraybackslash}p{2.4cm}|R|}
        \hline
        \ras{المرحلة} & \ras{دورها} \\ \hline
        \textbf{البيانات} & أمثلة كثيرة: نصوص أو صور أو أصوات. \\ \hline
        \textbf{التدريب}  & يتعلم النظام من البيانات، ويكتشف العلاقات بينها. \\ \hline
        \textbf{النموذج}  & ما ينتج عن التدريب: برنامج «تعلّم» القاعدة. \\ \hline
        \textbf{الاختبار} & نجرب النموذج على أمثلة جديدة لم يرها من قبل. \\ \hline
        \textbf{النتيجة}  & توقع أو تصنيف أو إجابة يقدمها النموذج. \\ \hline
    \end{tabularx}
\end{frame}

\begin{frame}{بيانات التدريب وبيانات الاختبار}
    \begin{tbox}{\faLightbulb\quad تعريف}
        \textbf{بيانات التدريب:} البيانات التي يتعلم منها النموذج ويكتشف فيها الأنماط.\\
        \textbf{بيانات الاختبار:} بيانات جديدة نستعملها لتقييم النموذج بعد تدريبه.
    \end{tbox}
    نقسم البيانات عادة إلى جزأين: جزء كبير للتدريب، وجزء صغير للاختبار.
    \begin{center}
    \begin{tikzpicture}
        \fill[brandTeal, rounded corners=3pt] (0,0) rectangle (-9.6,0.75);
        \fill[brandOrange, rounded corners=3pt] (-9.6,0) rectangle (-12,0.75);
        \node[lawn=white, font=\khatt\bfseries] at (-4.8,0.375) {بيانات التدريب (80\%)};
        \node[lawn=white, font=\khatt\bfseries] at (-10.8,0.375) {الاختبار (20\%)};
    \end{tikzpicture}
    \end{center}
\end{frame}

\begin{frame}{لماذا لا نختبر بالأمثلة نفسها؟}
    \begin{obox}{\faExclamationTriangle\quad انتبه!}
        لو اختبر الأستاذ التلميذ بالجمل نفسها التي شرحها في القسم،
        فقد يجيب التلميذ من الذاكرة لا من الفهم.\\[3pt]
        كذلك النموذج: نختبره ببيانات جديدة لنعرف هل \textbf{فهم} حقًا.
    \end{obox}
\end{frame}

\begin{frame}{تنوع البيانات}
    يتعلم النموذج ممّا يراه فقط. إذا كانت البيانات ناقصة أو غير متوازنة، كان التعلم ناقصًا.\\
    \textbf{مثال:} 500 صورة لهواتف، و500 صورة لحواسيب، و100 صورة فقط لألواح رقمية.
    \begin{center}
    \begin{tikzpicture}[y=0.0055cm]
        \draw[thick, brandMuted] (0.5,0) -- (-8.5,0);
        \fill[brandTeal]   (0,0)  rectangle (-2,500);
        \fill[brandTeal]   (-3,0) rectangle (-5,500);
        \fill[brandOrange] (-6,0) rectangle (-8,100);
        \node[above, font=\bfseries] at (-1,500) {500};
        \node[above, font=\bfseries] at (-4,500) {500};
        \node[above, font=\bfseries] at (-7,100) {100};
        \node[below, font=\khatt] at (-1,0) {هواتف};
        \node[below, font=\khatt] at (-4,0) {حواسيب};
        \node[below, font=\khatt] at (-7,0) {ألواح رقمية};
    \end{tikzpicture}
    \end{center}
\end{frame}

\begin{frame}{البيانات الجيدة}
    \begin{obox}{\faExclamationTriangle\quad بيانات غير ممثِّلة}
        صور الألواح الرقمية قليلة جدًا،
        لذلك سيخطئ النموذج كثيرًا عندما يرى لوحًا رقميًا جديدًا.
    \end{obox}
    \begin{tbox}{\faCheckCircle\quad البيانات الجيدة}
        البيانات الجيدة \textbf{متنوعة} و\textbf{ممثِّلة}:
        تغطي مختلف الحالات التي سيواجهها النظام في الواقع.
    \end{tbox}
\end{frame}

\begin{frame}{الانحياز الخوارزمي}
    \begin{tbox}{\faLightbulb\quad تعريف}
        \textbf{الانحياز الخوارزمي} نتائج غير عادلة أو غير دقيقة،
        سببها مشكلة في البيانات أو في طريقة تصميم النظام.
    \end{tbox}
    \begin{center}
        \flow[3cm]{{بيانات ناقصة\\أو من فئة واحدة}, {نموذج تعلّم\\صورة ناقصة}, {نتائج منحازة\\وغير عادلة}}
    \end{center}
\end{frame}

\begin{frame}{مثال من قسم اللغة العربية}
    \small
    \begin{tbox}{\faChalkboardTeacher\quad نموذج لتصحيح الأخطاء}
        نموذج لتصحيح الأخطاء تدرّب على نصوص صحفية حديثة فقط.
        عندما نعطيه بيتًا من الشعر الجاهلي، قد يعدّ بعض الكلمات الفصيحة القديمة أخطاءً.\\[3pt]
        \textbf{السبب:} لم يرَ هذا النوع من النصوص أثناء التدريب.
    \end{tbox}
    \textbf{\color{tealDeep}كيف نقلل الانحياز؟}
    \begin{itemize}
        \item نجمع بيانات \textbf{متنوعة} تمثل جميع الفئات والحالات.
        \item نوازن بين عدد الأمثلة في كل فئة.
        \item نختبر النموذج على حالات مختلفة، ونراجع أخطاءه.
    \end{itemize}
\end{frame}

\begin{frame}{خلاصة الدرس}
    \begin{tbox}{\faCheckCircle\quad خلاصة الدرس}
        \begin{itemize}
            \item تتعلم الآلة من \textbf{البيانات}، كما يتعلم التلميذ من الأمثلة.
            \item المراحل: البيانات $\leftarrow$ التدريب $\leftarrow$ النموذج
                  $\leftarrow$ الاختبار $\leftarrow$ النتيجة.
            \item نختبر النموذج دائمًا ببيانات \textbf{جديدة} لم يرها.
            \item البيانات الناقصة تؤدي إلى \textbf{الانحياز}، فنحتاج بيانات متنوعة وممثِّلة.
            \item النموذج يعكس البيانات التي تعلّم منها، بما فيها من \textbf{نقص}.
        \end{itemize}
    \end{tbox}
\end{frame}

\end{document}
